%!TEX root = ../main.tex

\section{Программа для задачи в полной постановке}

\subsection{Общая схема вычислений}

Научным коллективом разрабатывается программа для моделирования описанной задачи \eqref{eq:electrostatic}, \eqref{eq:phi} в полной, трехмерной, постановке. Программа использует параллельные вычисления с помощью MPI (запланирована также поддержка CUDA) и предназначена для работы на суперкомпьютерных комплексах.

Кратко обозначим схему вычислений. Приведенное описание во многом следует работе \cite{zipunova_computational}.

Расчетная область $\clOmega = [0, L]_x \times [0, L]_y \hm \times [0, L]_z$ имеет форму куба; в области $\clOmega$ вводится равномерная ортогональная расчетная сетка $\omega_h$ с шагом $h$ по каждой из координатных осей, $M = L / h \in \Natural$ -- число шагов по пространству. По времени сетка $\omega_h$ также равномерная, с шагом $\tau$. Целые пространственные индексы $j_{1,2,3} = \overline{0, M - 1}$ сопоставлены центрам ячеек сетки $\omega_h$; для обозначения точки на сетке используется три нижних, пространственных, индекса и один верхний, временной. Таким образом, $\omega_{j_1,j_2,j_3}^k = [(j_1 + 1/2)h \vect{e}_x + (j_2 + 1/2)h \vect{e}_y + (j_3 + 1/2)h \vect{e}_z, k \tau]$.

Свободным переменным $\phi$, $\Phi$, $\rho$ модели соответствуют сеточные функции; их значения отнесены к центрам ячеек сетки $\omega_h$. В качестве аргументов функции принимают описанный выше набор целых индексов: $\phi_{j_1,j_2,j_3}^k$ соответствует значению $\phi(\omega_{j_1,j_2,j_3}^k)$. 

Вводятся также двойственные сетки $\omega^*_\alpha$, $\alpha = 1,2,3$. Каждая $\omega^*_\alpha$ соответствует центрам граней исходной сетки $\omega_h$, перпендикулярных пространственной координатной оси $\mathcal{O}x_\alpha$. У сетки $\omega^*_\alpha$ индекс, соответствующий этой оси, является полуцелым: $j_\alpha \hm = -1/2 + \overline{0, M}$; прочие индексы совпадают с $\omega_h$. Координаты точки на сетке для $\omega^*_\alpha$ вычисляются по той же формуле $\omega_{j_1,j_2,j_3}^k$, приведенной выше.

Далее в случае, если некоторый координатный индекс не меняется в рамках целой формулы, он будет опускаться.

Пусть $H(\omega)$ -- пространство функций, заданных на сетке $\omega$. Введем разностные операторы $\delta_\alpha: H(\omega_h) \to H(\omega^*_\alpha)$; $\delta^*_\alpha: H(\omega^*_\alpha) \to H(\omega_h)$; $\mathring{\delta}_\alpha: H(\omega_h) \to H(\omega_h)$. Считая $j_\alpha$ целым,
\[
    [\delta_\alpha \zeta]_{j_\alpha + 1/2} = \cfrac{\zeta_{j_\alpha + 1} - \zeta_{j_\alpha}}{h}, \quad [\delta^*_\alpha \eta]_{j_\alpha} = \cfrac{\eta_{j_\alpha + 1/2} - \eta_{j_\alpha - 1/2}}{h}, \quad [\delta^*_\alpha \zeta]_{j_\alpha} = \cfrac{\zeta_{j_\alpha + 1} - \zeta_{j_\alpha - 1}}{2h}.
\]
Заметим, что $\mathring{\delta}_\alpha = \delta^*_\alpha \circ \delta_\alpha$. Все три оператора -- секточные приближения дифференцирования $\partial / \partial x_\alpha$ вдоль оси $\mathcal{O}x_\alpha$.

Также в ряде случаев при построении разностной схемы будет нужно переводить функцию с сетки $\omega_h$ на $\omega^*_\alpha$. Для краткости обозначений введем операторы
\[
    [s_\alpha \zeta]_{j_\alpha + 1/2} = \cfrac{\zeta_{j_\alpha + 1} - \zeta_{j_\alpha}}{2}, \qquad [g_\alpha \zeta]_{j_\alpha + 1/2} = \left( \cfrac{2}{1 / \zeta_{j_\alpha + 1} + 1 / \zeta_{j_\alpha}} \right)^{-1}.
\]
Функции на сетке $\omega^*_\alpha$ получаются осреднением значений исходной функции: в смысле среднего арифметического ($s_\alpha$) либо геометрического ($g_\alpha$).

\begin{subequations}
\label{sch:complete}
Перейдем к описанию схемы расчета. На каждом временном слое вначале решается разностный аналог уравнения \eqref{eq:electrostatic_c}
\begin{equation}
    \Div_h (\epsilon[\phi^k] \cdot \nabla_h \Phi^{k + 1}) = -\rho^k
    \label{sch:complete_a}
\end{equation}
относительно нового значения электрического потенциала $\Phi^{k + 1}$. Здесь $\Div_h$ и $\nabla_h$ обозначены сеточные приближения дифференциальных операторов, конкретный вид которых сейчас неважен. Решение уравнения происходит итерационно, многосеточным методом.

Далее делается шаг явной схемы для функций $\rho$ и $\phi$ по разностным аппроксимациям уравнений \eqref{eq:electrostatic_b} и \eqref{eq:phi}. Производную $\partial / \partial t$ в обоих случаях приблизим простейшей разностной производной первого порядка. Оператору $\Div$ везде будет соответствовать сумма $\delta^*_\alpha$; градиенту под знаком $\Div$ -- $\delta_\alpha$. $\sigma[\phi^k]$ в уравнении \eqref{eq:electrostatic_b} необходимо перевести на $\omega^*_\alpha$ -- добьемся этого взятием среднего геометрического $g_\alpha$. Операции $(\nabla \zeta, \nabla \zeta) \equiv \| \nabla \zeta \|_2^2$ сопоставим
\[
    \| \zeta \|_h^2 = \sum\limits_{\alpha = 1,2,3} (\mathring{\delta}_\alpha \zeta)^2.
\]
В уравнении \eqref{eq:phi} она встречается дважды: для $\Phi$, где описанной конструкции достаточно, и для $\phi$ под знаком $\Div$, где дополнительно требуется перевод функции на $\omega^*_\alpha$, что сделаем оператором среднего арифметического $s_\alpha$. Таким образом, получим разностные уравнения
\begin{gather}
    \cfrac{\rho^{k + 1} - \rho^k}{\tau} = \sum\limits_{\alpha = 1,2,3} \delta^*_\alpha(g_\alpha \sigma[\phi^k] \cdot \delta_\alpha \Phi^{k + 1}),
    \label{sch:complete_b} \\
    \cfrac{\phi^{k + 1} - \phi^k}{\tau} = \half \epsilon'(\phi^k) \| \nabla \Phi^{k + 1} \|_h^2 + \cfrac{\Gamma}{l^2} f'(\phi^k) + \sum\limits_{\alpha = 1,2,3} \delta^*_\alpha \left( \cfrac{\Gamma}{2} \delta_\alpha \phi^k + \beta \Gamma l^2 \cdot s_\alpha \| \nabla \phi^k \|_h^2 \cdot \delta_\alpha \phi^k \right),
    \label{sch:complete_c}
\end{gather}
откуда легко выражаются $\rho^{k + 1}$ и $\phi^{k + 1}$.
\end{subequations}

На границе области $\clOmega$ заданы граничные условия функций $\phi$ и $\Phi$ типа Дирихле либо однородные типа Неймана, согласно \eqref{cond:boundary_basic}. Считаем, что $\Phi|_{y = L} = \Phi^- = 0$; $\Phi|_{y = 0} \hm = \Phi^+$ в текущей постановке задачи принимается равным константе (условие \eqref{cond:boundary_voltage}, типа Дирихле), однако может рассматриваться и как числовая свободная переменная модели (условие \eqref{cond:boundary_balance} или \eqref{cond:boundary_charge}).

Для аппроксимации граничных условий на каждой грани $\clOmega$ добавим к сетке $\omega$ по одному слою фиктивных ячеек (с индексами $-1$, $M$ по соответствующей координате). Теперь разностная производная $\delta_\alpha$ определена на гранях с $j_\alpha = -1/2, M - 1/2$; производная $\mathring{\delta}_\alpha$ -- в ячейках с $j_\alpha = 0, M - 1$. Для условия Дирихле $\phi|_{x_\alpha = 0, L} = \psi$ примем $\phi^k_{j_\alpha = -1,M} = \psi(\omega_{j_\alpha = -1,M})$; для однородного условия Неймана перед шагом явной схемы проведем присваивание $\phi^k_{j_\alpha = -1,M} := \phi^k_{j_\alpha = 0, M - 1}$, так чтобы $[\delta_\alpha \phi]_{j_\alpha = -1/2, M - 1/2}^k = 0$.


\subsection{Адаптация шага по времени}

Раздел \ref{sec:adaptation} посвящен адаптации расчетного шага по времени в исследуемой модели, которая тогда была проведена для упрощенной, одномерной, постановки задачи. Однако выбранные методы не зависят от размерности, и основной целью изначально было внедрение адаптации в программу для моделирования полной постановки задачи.

В описанную программу, реализующую вычислительную схему \eqref{sch:complete}, введен метод адаптации расчетного шага по времени по фазовому полю, наиболее простой из трех рассмотренных в разделе \ref{sec:adaptation} и показавший в целом лучшие результаты в программе для одномерного случая задачи.

В уравнениях \eqref{sch:complete_b}, \eqref{sch:complete_c} шаг по времени $\tau$ заменяется на переменный $\tau^k$, рассчитываемый по формулам \eqref{sch:time_step_min_max}, \eqref{sch:time_step_phi}. В качестве $[\partial \phi / \partial t]_h$ в знаменателе выражения \eqref{sch:time_step_phi} берется значение левой части формулы \eqref{sch:complete_c} -- так как в расчете уже используется ее правая часть, дополнительных вычислений это не потребует. Однако из-за распределения счета между вычислительными узлами MPI использовать при переходе между временными слоями $k$ и $k + 1$ текущее значение производной затратно: необходима лишняя синхронизация процессов с массовой операцией максимума. Поэтому будет использоваться значение производной с предыдущего перехода, а именно:
\[
    \widetilde{\tau}^k = \tol \cdot \left( \max\limits_{j_{1,2,3}} \left| \cfrac{\phi^k - \phi^{k - 1}}{\tau^{k - 1}} \right|_{j_{1,2,3}} \right)^{-1}.
\]